\begin{tomtat}
\begin{itemize}
  \item Ba lối thoát khác nhau ở chỗ chúng làm gì với vế phải của
định nghĩa~\ref{def:ch01-do-trung-thuc}, tức cái mà hành vi mô hình phụ thuộc
vào và không ai biết với một mô hình có sẵn. Đường thứ nhất ép nó phải biết
trước, đường thứ hai thay nó bằng đánh giá của người đọc, đường thứ ba mở mô
hình ra đọc nó. Chữ \enquote{lối thoát} là chữ của sách: trong bốn bài của
chương, chỉ bài của đường thứ nhất tự đứng ở chỗ ấy, và chữ \enquote{faithful}
không xuất hiện lần nào trong phần thân của ba bài kia. Riêng đường thứ nhất có
bốn dạng chứ không phải một: thiết kế nhiệm vụ, thiết kế dữ liệu, thiết kế mô
hình, và nhãn khái niệm có sẵn trong bộ dữ liệu. Ba dạng đầu ép một thứ vào chỗ
nó vốn không có, nên kết luận rút ra nói về thiết lập đã ép; dạng thứ tư không ép
gì, và đổi lại đòi một điều kiện có trước là bộ dữ liệu phải đã được gán nhãn và
tập nhãn ấy phải đủ dự báo nhiệm vụ.
  \item Cái giá của đường thứ nhất đo được trên chính bảng của bài báo: phần
nhiệm vụ do bài tự sinh cấp 486 trong 1946 nhãn mức bước và cấp đúng một trong
bốn loại bước mà bảng ấy đếm. Mục hạn chế của bài chỉ nêu hai chỗ, phủ sóng
trong một chuỗi và lệch nhãn, nên câu nói rằng đường này mua tiêu chí bằng tính
đại diện là câu
của sách, và bài trả lời trước lời buộc tội ấy bằng cách ghi rằng nhiệm vụ được
dựng còn chuỗi suy luận thì thật.
  \item Đường thứ hai xác nhận được lời đề xuất rằng một chỉ số thay cho chất
lượng người dùng cảm nhận, và không xác nhận được lời đề xuất rằng một chỉ số đo
độ trung thực, vì người đọc không nhìn thấy bên trong mô hình. Nó đã được chạy
128 lần trên 18\,254 bài. Cái giá của nó không phải tiền: bài in các chữ
\enquote{expensive} và \enquote{power} không lần nào và chữ \enquote{cost} đúng
một lần trong một tên riêng, nêu hai lời phản đối thông thường rồi bác cả hai
bằng câu rằng những khó khăn ấy không phải là không vượt qua nổi, và cái nó ghi lại là chất lượng của bằng
chứng, với 13 trong 128 bài không báo số người tham gia.
  \item Đường thứ ba, trong bài duy nhất vẽ bản đồ cả ba nhánh attribution, là
nhánh không có ground truth nào được nêu. Bài ấy xếp
\tn{diễn giải cơ chế}{mechanistic interpretability} vào bản đồ chứ không làm nó,
khung hình thức của nó chỉ phủ attribution đặc trưng và việc mở rộng là một giả
thuyết, và nó phát biểu chỗ thiếu dụng cụ đo cho cả ba nhánh cùng lúc. Thủ tục
của nhánh này còn đòi hai thứ có trước, một bộ dữ liệu gợi ra đúng hành vi cần
xét và một chỉ số nói được rằng đã tìm thấy, tức hai lựa chọn tự do của
mục~\ref{sec:ch08-vong-lap-chung} chuyển từ đặc trưng sang thành phần.
  \item Một cái tên đọc được nghĩa là một công cụ dò chưa ai kiểm đã gán cho một
nơ-ron cái tên khớp hơn mức trung bình của chính từ điển ấy, và nghĩa hết ở đó.
Tỉ lệ nơ-ron đạt cả hai tính chất chạy từ 11.88 tới 19.82 phần trăm trên ba mô
hình phía sau, thành 29.14 phần trăm nếu đổi mốc sang điểm trung bình của mô
hình gốc, và điểm đọc được giãn gấp khoảng 1.43 lần theo công cụ chấm. Ở một tập
khác, 1329 khái niệm của tầng thắt bù, điểm điều khiển được tương quan
$r = 0.06$ giữa hai mô hình phía sau, mà chính bài giải thích bằng việc hai mô
hình ấy không dùng cùng bộ mã hóa.
  \item Ba đường không song song, và đó là câu của sách. Đường thứ nhất được
chọn vì bài của nó xét đường thứ ba là không dùng được, đường thứ ba thừa kế chỗ
thiếu mà nó lẽ ra phải lấp, và đường thứ hai bị bỏ không phải vì đắt. Một đường
thu hẹp câu hỏi, một đường đổi câu hỏi, một đường hoãn câu hỏi.
\end{itemize}
\end{tomtat}

\cauhoi
\begin{cauhoids}
  \item Mục~\ref{sec:ch16-gia-duong-dung-san} cho thấy với ngưỡng 1 và với ngưỡng
4 trên hàm chạy suốt sách, đúng một trong hai câu trả lời chỉ ra được tập phụ
thuộc. Hãy cho biết có ngưỡng nào mà \emph{cả hai} câu trả lời đều chỉ ra được
tập phụ thuộc hay không. Nếu có thì nêu ngưỡng ấy và viết ra tập phụ thuộc ở cả
bốn tổ hợp; nếu không thì giải thích vì sao không.
  \item Cột nhiệm vụ tự sinh cấp 486 trong 1946 nhãn mức bước và cấp đúng loại
bước thực thi bước nút thắt. Hãy phát biểu một câu về phương pháp dựng ground
truth theo thiết kế nhiệm vụ mà con số ấy ủng hộ, rồi phát biểu một câu nghe
giống nó mà con số ấy \emph{không} ủng hộ, và nói chỗ khác nhau giữa hai câu nằm
ở đâu.
  \item Mục~\ref{sec:ch16-duong-thu-hai} ghi rằng bài của đường thứ hai bác lời
giải thích theo chi phí. Hãy nêu một quan sát thực nghiệm mà nếu đo được thì nó
sẽ ủng hộ lời giải thích theo chi phí, và một quan sát khác ủng hộ lời giải
thích theo cấu trúc khuyến khích mà bài đưa ra, rồi cho biết cuộc kiểm đếm của
bài có phân biệt được hai lời giải thích ấy hay không.
  \item Xếp cả ba lối thoát vào bảng~\ref{tab:ch01-bon-quan-he}: mỗi đường cấp
bằng chứng cho dòng nào. Với đường nào không xếp gọn vào một dòng thì nói rõ nó
cấp bằng chứng cho một phần của dòng nào và không cấp cho phần nào.
  \item Mục~\ref{sec:ch16-duong-thu-ba} đọc yêu cầu về bộ dữ liệu gợi ra hành vi
và về chỉ số đo thành công như hai lựa chọn tự do của
mục~\ref{sec:ch08-vong-lap-chung} chuyển sang mức thành phần. Hãy nêu lựa chọn
tự do thứ ba của mục ấy và cho biết ở mức thành phần thì nó tương ứng với cái
gì, hoặc vì sao nó không tương ứng với cái gì cả.
  \item Tỉ lệ nơ-ron đạt cả hai tính chất là 18.84 phần trăm dưới mốc trung bình
của chính từ điển và 29.14 phần trăm dưới mốc trung bình của mô hình gốc. Hãy
cho biết mỗi mốc trả lời câu hỏi nào, rồi chọn mốc mà bạn cho là hợp với câu hỏi
của chương này hơn và bảo vệ lựa chọn ấy bằng một câu.
  \item Hệ số $r = 0.06$ đo trên 1329 khái niệm của tầng thắt bù ở hai mô hình
phía sau, và bài giải thích nó bằng việc hai mô hình ấy không dùng cùng bộ mã
hóa. Hãy cho biết lời giải thích ấy làm yếu đi kết luận nào trong ba kết luận
sau và để nguyên kết luận nào: rằng phép đo khả năng điều khiển không đáng tin,
rằng khả năng điều khiển là tính chất của một cặp chứ không của một khái niệm,
và rằng một khái niệm điều khiển được ở mô hình này thì cũng điều khiển được ở
mô hình kia. Sau đó nêu một thí nghiệm tách được nguyên nhân mà bài nêu ra khỏi
những nguyên nhân khác.
  \item Đọc phụ lục F.3, \enquote{Towards fairer and more practical evaluation},
trong \enquote{Towards Unified Attribution in Explainable AI, Data-Centric AI,
and Mechanistic Interpretability} của Zhang và cộng sự~\autocite{p18unified}.
Phụ lục ấy nêu ba cách đánh giá và một chỗ hỏng cho mỗi cách.
Mục~\ref{sec:ch16-duong-thu-ba} đã chỉ ra rằng chỉ hai trong ba cách ấy tương
ứng với hai lối thoát của chương này. Hãy nêu cách thứ ba, cho biết nó đo cái gì,
và giải thích vì sao nó không phải một lối thoát theo nghĩa
mục~\ref{sec:ch16-ba-loi-thoat} đặt cho chữ ấy.
\end{cauhoids}
